\section*{Chapter \ref{ch:b2:chemical-potential} --- Chemical Potential and Mixtures}

\begin{solution}{exo:b2:chemical-potential:1}
The tangent at $x_2 = 0.4$ cuts the edges at $\bar V_1 \approx 14$ and $\bar
V_2 \approx 49$ (model units). Then $0.6 \times 14 + 0.4 \times 49 = 28.0$,
the molar volume read on the curve at $x_2 = 0.4$, below the
$0.6 \times 18 + 0.4 \times 58 = 34.0$ of the pure volumes.
\end{solution}

\begin{solution}{exo:b2:chemical-potential:2}
$n(\text{water}) = 1000/18.015 = \qty{55.5}{mol}$, $x_1 = 55.5/56.5 = 0.982$;
$p = 0.982 \times 3.17 = \qty{3.11}{kPa}$: a lowering of 1.8\,\%.
\end{solution}

\begin{solution}{exo:b2:chemical-potential:3}
$p(\ce{O2}) = 0.2095 \times \qty{1.013e5}{Pa} = \qty{2.12e4}{Pa}$;
$c = 1.3 \times 10^{-5} \times 2.12 \times 10^4 = \qty{0.28}{mol/m^3} =
\qty{2.8e-4}{mol/L}$, that is $2.8 \times 10^{-4} \times 32.0 \times 10^3 =
\qty{8.8}{mg/L}$.
\end{solution}

\begin{solution}{exo:b2:chemical-potential:4}
$Q = (1.0)^2/[(1.0)(0.010)^3] = 1.0 \times 10^6$; $RT\ln Q = 2.479 \times 13.82 =
\qty{34.2}{kJ/mol}$; $\Delta_r G = -32.8 + 34.2 = \qty{+1.4}{kJ/mol} > 0$: in this
mixture, so poor in hydrogen, ammonia decomposes.
\end{solution}

\begin{solution}{exo:b2:chemical-potential:5}
Gibbs--Duhem: $x_1\dd\bar V_1 + x_2\dd\bar V_2 = 0$, so $\dd\bar V_2 =
-(0.70/0.30) \times (-0.30) = \qty{+0.70}{mL/mol}$.
\end{solution}

\begin{solution}{exo:b2:chemical-potential:6}
$c(\text{particles}) = 2 \times 9.0/58.44 = \qty{0.308}{mol/L} =
\qty{308}{mol/m^3}$; $\Pi = 308 \times 8.314 \times 310 = \qty{7.9e5}{Pa}$,
about \qty{7.9}{bar}. The inside of a red cell has the same \omterm{def:b2:chemical-potential:osmotic-pressure}{osmotic pressure}:
no net flow of water. In pure water, water enters the cell, which swells and
bursts.
\end{solution}

\begin{solution}{exo:b2:chemical-potential:7}
$c = \Pi/RT = 825/(8.314 \times 298.15) = \qty{0.333}{mol/m^3}$; in
\qty{1.000e-4}{m^3}, $n = \qty{3.33e-5}{mol}$, $M = 2.00/3.33 \times 10^{-5}
= \qty{6.0e4}{g/mol}$ (\qty{60}{kg/mol}). $\Delta T_f = 1.86 \times 3.33
\times 10^{-5}/0.100 = \qty{6.2e-4}{K}$: unmeasurable, while the \omterm{def:b2:chemical-potential:osmotic-pressure}{osmotic
pressure} is a column of \qty{8}{cm} of water.
\end{solution}

\begin{solution}{exo:b2:chemical-potential:8}
(a) $I = \qty{1.0e-3}{mol/kg}$, $\log\gamma_\pm = -0.51 \times 0.0316$,
$\gamma_\pm = 0.964$. (b) $I = \frac12(4 \times 1.0 + 1 \times 2.0) \times
10^{-3} = \qty{3.0e-3}{mol/kg}$, $\log\gamma_\pm = -0.51 \times 2 \times 0.0548$,
$\gamma_\pm = 0.879$. (c) $I = \frac12(4 + 4) \times 10^{-3} = \qty{4.0e-3}{mol/kg}$,
$\log\gamma_\pm = -0.51 \times 4 \times 0.0632$, $\gamma_\pm = 0.74$.
\end{solution}

\begin{solution}{exo:b2:chemical-potential:9}
$b = 0.310/1.86 = \qty{0.167}{mol/kg}$; $n = 0.167 \times 0.0500 =
\qty{8.33e-3}{mol}$; $M = 1.50/8.33 \times 10^{-3} = \qty{180}{g/mol}$ (a
hexose sugar, for instance).
\end{solution}

\begin{solution}{exo:b2:chemical-potential:10}
$x_1\,\dd(Ax_2^2) + x_2\,\dd(Ax_1^2) = 2Ax_1x_2(\dd x_2 + \dd x_1) = 0$ since
$x_1 + x_2 = 1$. As $x_1 \to 0$, $\gamma_1 \to \mathrm e^{A}$ and $p_1 =
\gamma_1x_1p_1^* \approx \mathrm e^{A}p_1^*\,x_1$: $k_{H,1} = p_1^*\mathrm e^A$.
For $A = 1.1$, $\mathrm e^{1.1} = 3.0$: Henry's slope is three times Raoult's.
\end{solution}

\begin{solution}{exo:b2:chemical-potential:11}
(a) A volatile solute contributes its own vapour pressure; the vapour is no
longer pure solvent and the derivation (pure vapour of the solvent in
equilibrium with the solution) fails. (b) If the solute enters the solid, the
solid's \omterm{def:b2:chemical-potential:chemical-potential}{chemical potential} is lowered too, by $RT\ln x_1^{\text{s}}$. If the
solid solution has the same composition as the liquid, both lines move down
by the same amount and the freezing point does not change: the depression
needs a solute rejected by the crystal.
\end{solution}

\begin{solution}{exo:b2:chemical-potential:12}
Cryoscopy: $b_{\min} = 0.01/1.86 = \qty{5.4e-3}{mol/kg}$; ebullioscopy:
$0.01/0.513 = \qty{1.9e-2}{mol/kg}$. At $\qty{5.4e-3}{mol/L} =
\qty{5.4}{mol/m^3}$, $\Pi = 5.4 \times 8.314 \times 298 = \qty{1.3e4}{Pa}$,
more than a metre of water; an osmometer reads a few pascals. Osmometry is by
far the most sensitive, which is why it measures the large molar masses of
polymers, whose molalities are tiny.
\end{solution}

\begin{solution}{pb:b2:chemical-potential:1}
\textbf{1.} $\mu_1 = \mu_1^*(T, p) + RT\ln x_1$.
\textbf{2.} Diol $30/62.07 = \qty{0.483}{mol}$, water $70/18.015 =
\qty{3.886}{mol}$; $x_1 = 3.886/4.369 = 0.889$.
\textbf{3.} $p = 0.889 \times 3.17 = \qty{2.82}{kPa}$.
\textbf{4.} $b = 0.483/0.070 = \qty{6.90}{mol/kg}$.
\textbf{5.} The diol boils at \qty{197}{\celsius}: its vapour pressure at room
temperature is far below that of water.
\textbf{6.} $\mu_1^{\text{s}}(T_f) = \mu_1^{*\text{l}}(T_f) + RT_f\ln x_1$.
\textbf{7.} $\ln x_1 = -\Delta_{\text{fus}}G(T)/(RT)$ with
$\Delta_{\text{fus}}G = \mu_1^{*\text{l}} - \mu_1^{\text{s}}$, zero at $T_f^*$;
Gibbs--Helmholtz, $\dd(\Delta_{\text{fus}}G/T)/\dd T = -\Delta_{\text{fus}}H/T^2$,
integrated from $T_f^*$ to $T_f$, gives $\Delta_{\text{fus}}G(T_f)/T_f =
\Delta_{\text{fus}}H(1/T_f - 1/T_f^*)$, whence the relation.
\textbf{8.} $\ln x_1 \approx -x_2 \approx -bM_1$ and $1/T_f - 1/T_f^* \approx
-\Delta T_f/T_f^{*2}$: $\Delta T_f = (RT_f^{*2}M_1/\Delta_{\text{fus}}H)\,b$.
\textbf{9.} $333.4 \times 18.015 = \qty{6007}{J/mol}$; $K_f = 8.314 \times
273.15^2 \times 0.018015/6007 = \qty{1.86}{K.kg/mol}$.
\textbf{10.} That $\Delta_{\text{fus}}H$ does not change between $T_f$ and
$T_f^*$ (it does, by Kirchhoff, about 2\,\% per \qty{3}{K} here).
\textbf{11.} $\Delta T_f = 1.86 \times 6.90 = \qty{12.8}{K}$: \qty{-12.8}{\celsius}.
\textbf{12.} $1/T_f = 1/273.15 - (8.314/6007)\ln 0.889$: $T_f = \qty{261.6}{K}$,
\qty{-11.6}{\celsius}.
\textbf{13.} The exact ideal relation, which does not assume a dilute
solution (here 11\,\% of the molecules are diol). Both assume an \omterm{def:b2:chemical-potential:ideal-mixture}{ideal mixture}
and a constant enthalpy of fusion; water and diol, linked by hydrogen bonds,
are not ideal, so the real freezing point differs.
\textbf{14.} Pure ice: the diol stays in the liquid.
\textbf{15.} Removing water as ice lowers $x_1$ in the liquid, and by the
relation the equilibrium temperature falls: the coolant freezes over a range
of temperatures, not at one point.
\textbf{16.} $K_b = 8.314 \times 373.12^2 \times 0.018015/40\,650 =
\qty{0.513}{K.kg/mol}$.
\textbf{17.} $\Delta T_b = 0.513 \times 6.90 = \qty{3.5}{K}$.
\textbf{18.} $1/T = 1/373.12 - (8.314/40\,650)\ln(2.0/1.013)$: $T = \qty{393.5}{K}$,
about \qty{120}{\celsius}.
\textbf{19.} The cap: \qty{20}{K} against \qty{3.5}{K} for the diol.
\textbf{20.} $x_1 = \exp[-(6007/8.314)(1/253.15 - 1/273.15)] = 0.811$.
\textbf{21.} $x_2 = 0.189$: $w = 0.189 \times 62.07/(0.189 \times 62.07 + 0.811
\times 18.015) = 0.44$.
\textbf{22.} Dilute law: $b = 20/1.86 = \qty{10.75}{mol/kg}$, $w = 10.75 \times
62.07/(1000 + 10.75 \times 62.07) = 0.40$. The linearisations ($\ln x_1 \approx
-x_2$, $x_2 \approx bM_1$) overstate the lowering per mole of diol at such a
high concentration, so the dilute law asks for less diol.
\textbf{23.} In the ideal-mixture model, a coolant with a mass fraction of
diol $\boldsymbol{w \approx 0.44}$ stays liquid down to \qty{-20}{\celsius}.
\end{solution}
